\section{A projekt kiinduló problémája}
Manapság az emberek filmnézési szokásai több platformra terjednek ki: streaming szolgáltatókra, mozikra, fizikai adathordozókra és televíziós tartalmakra. Ezek a megoldások rendkívül kézenfekvővé teszik a filmek fogyasztását, ugyanakkor egy új problémát is létrehoznak. Egy adott szolgáltatás, például a Netflix, képes nyilvántartani, hogy a felhasználó az adott platformon milyen filmeket nézett meg, vagy mit szeretne később megnézni, ez azonban nem terjed ki más szolgáltatásokra.

Emiatt a felhasználó filmnézési előzményei nehezen követhetővé válnak. Előfordulhat, hogy valaki nem emlékszik pontosan, látott-e már egy filmet, milyen értékelést adott volna rá, vagy mely filmeket szerette volna később megnézni. A probléma tehát nem kizárólag a filmek megtalálásáról szól, hanem egy személyes, platformfüggetlen filmnapló kialakításáról is.

A filmnézésnek ugyanakkor közösségi oldala is van. Gyakran kíváncsiak vagyunk arra, hogy ismerőseink milyen filmeket néztek meg az utóbbi időben, mit gondoltak róluk, vagy milyen filmeket ajánlanának. Ezek az ajánlások sokszor beszélgetésekben, csoportos üzenetekben vagy közösségi médiás bejegyzésekben jelennek meg, így könnyen elvesznek, és később nehezen kereshetők vissza. A legtöbb streaming platform elsősorban a saját katalógusára és az egyéni fogyasztásra koncentrál, ezért nem ad teljes választ arra, hogyan lehetne a személyes filmnaplót és a baráti ajánlásokat egységes rendszerben kezelni.

A WatchDB projekt kiinduló problémája ezért kettős. Egyrészt szükség van egy egyszerű, platformfüggetlen felületre, ahol a felhasználó nyilvántarthatja a megnézett filmeket és a később megtekintendő alkotásokat. Másrészt megjelenik az igény egy olyan közösségi rétegre is, amelyen keresztül a felhasználók követhetik egymás filmnézési aktivitását és ajánlásokat fedezhetnek fel. A dolgozatban bemutatott projekt ezt a problémát használja kiindulópontként, majd ebből vezeti le a termék koncepcióját, a követelményeket és az MVP-hez szükséges funkciókat.
